\section{Introduction}\label{sec:intro}
Constraint optimization is a general model for centralized and distributed problem solving, which has a wide range of applications in artificial intelligence and multiagent systems \cite{YedidsionZF18,BoothB19,GolestanianBTB23}.
Among the applications that can be represented as constraint optimization problems (COPs) are meeting scheduling, task allocation, bus driver scheduling, disaster response, and operating room assignment \cite{csplib,NelkeOZ20,LavieCLZ23,GaonSZ23}.
Two main approaches are commonly used to solve COPs. The first is search, in which a candidate solution is maintained and updated while traversing the solution space to find one with higher quality than those found previously \cite{ModiSTY05,Zivan08,YeohFK10,LeeMS15}.
	The second is dynamic programming inference \cite{PetcuF05,RogersFSJ11,CohenGZ20} in which information is aggregated so that it enables the selection of high-quality solutions. These inference algorithms are rooted in techniques such as belief propagation and bucket elimination \cite{Pearl88,Dechter99}, which are used to solve optimization problems represented by graphical models.

	Min-sum (also called Max-sum when applied to maximization problems) is an incomplete inference algorithm that has drawn considerable attention in recent years \cite{FarinelliRPJ08,YuCHD17,CerquidesEPR18,DengA20}. It is designed as a message passing algorithm in which nodes in a factor graph exchange messages with neighboring nodes. Thus, it was found to be useful for solving multi-agent applications such as sensor systems \cite{TeacyFGPRJ08,StrandersFRJ09,YedidsionZF18}, task allocation for rescue teams in disaster areas \cite{RamchurnFMJ10}, and IoT applications \cite{RustPR16}.  Min-sum is actually a version of the well known belief propagation algorithm \cite{Pearl88,YanoverMW06}, adjusted to solve COPs and distributed COPs (DCOPs) \cite{BritoM10,GhoshKV15,KhanTYJ18,CohenLZ23}.

	Belief propagation in general (and Min-sum specifically) is known to converge to the optimal solution when solving problems whose constraint graph is acyclic. For problems that include cycles, there is no such guarantee \cite{YanoverMW06,FarinelliRPJ08}.  Furthermore, in graphs with multiple cycles, duplicated information is propagated, leading to inaccurate and inconsistent inference \cite{Pearl88}.  Unfortunately, the underlying problem representation graphs of many realistic applications indeed include multiple cycles (e.g., \citet{ModiSTY05,GershmanMZ09}). Damping is a method that improves Min-sum's performance on problems with multiple cycles~\cite{CohenGZ20,ZaedLZ25}. However, while it triggers convergence of the algorithm to high quality solutions, thousands of iterations are required to reach these, and thus, damped Min-sum (DMS) is not a good choice when very fast decision making is required.
In \citet{CohenGZ20}, the splitting of nodes in the factor-graph upon which the algorithm performs, which represent functions (i.e., function-nodes), was found to trigger rapid convergence of DMS. The DMS algorithm was performed on a new generated factor-graph in which each of the function-nodes in the original factor-graph was represented by two function-nodes. The costs for combination of assignments in the original function-nodes were divided between the two function-nodes in the new generated factor-graph. This caused rapid convergence to high quality solutions. However, although the success of this version of the algorithm (DMS-SCFG) was demonstrated empirically, a formal theoretic explanation of this phenomenon was not provided.

In this study, we investigate the success of function-node splitting in triggering convergence of Min-sum in general and of DMS in particular. We identify a phenomenon that, to the best of our knowledge, has not been reported before: the small cycles generated by the splitting trigger a feedback loop that rapidly separates the solution to which the algorithm converges from alternative candidates. We establish explicit bounds within which the assignments of the variables cannot change, i.e., within which the algorithm converges, and we prove that this convergence survives bounded, dynamically changing influence from the rest of the factor graph.

We further reveal why the use of splitting without damping commonly results in oscillations between solutions of low quality. We demonstrate that Min-sum, when performed on a split constraint factor graph with no damping, tends to converge to two separate solutions between which the nodes of the graph alternate, unfortunately in an uncoordinated manner. This generates mixed assignments of low quality, between which the algorithm oscillates. This explains the importance of combining splitting with damping: damping enforces a strong relation between the calculations performed by nodes in consecutive iterations and thus prevents this phenomenon.

Finally, we use the theoretical knowledge we acquired in two ways. First, as a diagnostic: letting standard Min-sum with splitting reach the stage where it oscillates between two solutions and then selecting between the two with standard search recovers most of the quality that the uncoordinated oscillation hides, which confirms the two-solution reading but does not match the split-damped algorithm. Second, to design a heuristic: running DMS for a predetermined number of iterations and only then splitting the factor graph improves on immediate splitting on the random benchmarks, and is the best distributed variant on every benchmark family we tested.
